\documentclass[../../../main.tex]{subfiles}
\begin{document}
\subsection{Dirac Equation}
The Dirac equation is
\begin{equation*}
    i \hbar \; \partial_t \Psi(\mathbf{r },t) = H \Psi(\mathbf{r },t)
\end{equation*}
with the Hamiltonian
\begin{equation*}
    H = c \hat{\boldsymbol{\alpha}} \cdot \mathbf{p} + \hat{\beta}mc^2
    = \begin{bmatrix}
        m c^{2} \mathbf{1}                        & c \boldsymbol{\sigma}\cdot\hat{\mathbf{p}} \\
        c\boldsymbol{\sigma}\cdot\hat{\mathbf{p}} & -m c^{2} \mathbf{1}
    \end{bmatrix}
\end{equation*}
and $\Psi(\mathbf{r },t)$ as the four-component spinor.
It can also be written in this way
\begin{equation*}
    i\hbar \;\partial_t \Psi(\mathbf{r},t)=\left(-i\hbar c\hat{\boldsymbol{\alpha}}\cdot\nabla + \hat{\beta}m c^{2}\right)\Psi(\mathbf{r},t)
\end{equation*}

The matrices $\hat{\alpha}_i$ and $\hat{\beta}$ defined as
\begin{equation*}
    {\hat{\alpha}}_i = \begin{bmatrix} \mathbf{0} & {\sigma}_i\\  {\sigma}_i & \mathbf{0} \end{bmatrix}\qquad
    \beta = \begin{bmatrix} \mathbf{1} & \mathbf{0}\\ \mathbf{0} & -\mathbf{1} \end{bmatrix}
\end{equation*}
and satisfies
\begin{equation*}
    \{\hat{\alpha}_i,\alpha_j\}=2\delta_{ij} \qquad
    \{\hat{\alpha}_i,\beta\}=\mathbf{0} \qquad
    \hat{\alpha_i}=\hat{\beta}^2=\mathbf{1}
\end{equation*}

The equation has the same structure with the Schrödinger equation, in fact this is the relativistic generalization of the Schrödinger equation.
This equation is constructed specifically to satisfy the relativistic energy $E= \sqrt{\mathbf{p }^2c^2+m^2c^4}$.

\subsubsection{Spinor-field.}
The four-components spinor encode two states of spin-$1/2$ particle and degenerate solution of relativistic dispersion relation $E^2 = p^2c^2+m^2c^4$.
The components can be grouped as
\begin{equation*}
    \Psi = \begin{bmatrix}\phi\\\chi\end{bmatrix} \qquad
    \phi = \begin{bmatrix}\psi_1\\\psi_2\end{bmatrix} \qquad
    \chi = \begin{bmatrix}\psi_3\\\psi_4\end{bmatrix}
\end{equation*}
$\phi$ component encode degrees of freedom at low energy; while $\chi$ encode relativistic correction responsible for spin–orbit coupling, magnetic moment, relativistic effects.
The coupling of $\phi$ and $\chi$
\begin{equation*}
    \begin{bmatrix}
        m c^{2}                                   & c \boldsymbol{\sigma}\cdot\hat{\mathbf{p}} \\
        c\boldsymbol{\sigma}\cdot\hat{\mathbf{p}} & -m c^{2}
    \end{bmatrix}
    \begin{bmatrix}\phi\\\chi\end{bmatrix}
    = \begin{bmatrix}
        mc^{2}\phi + c\bigl(\boldsymbol{\sigma}\cdot\mathbf{p}\bigr)\chi \\
        c\bigl(\boldsymbol{\sigma}\cdot\mathbf{p}\bigr)\phi - mc^{2}
    \end{bmatrix}
\end{equation*}
via $\boldsymbol{\sigma}\cdot \mathbf{p}$  is what makes the theory relativistic.

\subsubsection{Klein-Gordon equation.}
The Klein-Gordon originally interpreted to generalize the Schrödinger equation in relativistic.
It starts with the relativistic energy
\begin{equation*}
    E^2 = p^2c^2+ m^2c^4
\end{equation*}
then introduce the prescription $    \mathbf{p} \to -i\hbar\nabla $ and  $    H \to i\hbar \;\partial_t$
\begin{equation*}
    \bigg(\frac{1}{c^2}\frac{\partial^2}{\partial t^2} - \nabla^2 + \frac{m^2 c^2}{\hbar^2}\bigg)\psi(\mathbf r,t) = 0
\end{equation*}
The equation is also known as the generalization of Maxwell’s wave equation for massive particles.
Aside, the equation suffer two problems.
\begin{enumerate}
    \item \textbf{Negative energy.}
          The equation was supposed to describe single relativistic particle $\psi(\mathbf{r},t)$, it however include both spectrum of $+E$ and $-E$, so the theory is unbounded below if interpreted as a single-particle system.

          In relativistic quantum field theory, this is corrected by interpreting $\psi(\mathbf{r},t)$ as the quantum field, and the negative energy as antiparticle.
    \item \textbf{Incorrect probability density.}
          A probability density must satisfy conserved
          \begin{equation*}
              \frac{d }{dt } \int d^3 \mathbf{r} \rho(\mathbf{r},t) = 0
          \end{equation*}
          and positive definite
          \begin{equation*}
              \rho(\mathbf{r},t) \geq  0
          \end{equation*}
          quantity.
          Single plane wave state $\rho(\mathbf{r},t) = |\psi|^2$ does produce time-independent $\rho(\mathbf{r})$, however superpositions of pure states cause interference term to appear and time dependence to be produced.
          There exists a conserved current
          \begin{equation*}
              \rho = \frac{i\hbar}{2 m c^{2}}\bigl(\psi^{*}\partial_{t}\psi - \psi\partial_{t}\psi^{*}\bigr)
          \end{equation*}
          that satisfies the time-independence.
          This probability density, however, is not positive definite, so it still invalidates its interpretation as a probability density.
\end{enumerate}

Dirac's equation able to produce the continuity equation
\begin{equation*}
    \partial_t \rho(\mathbf{r },t) + c \nabla \cdot \mathbf{j}(\mathbf{r },t)=0
\end{equation*}
with $\rho(\mathbf{r },t)$ and $\mathbf{j}(\mathbf{r },t)$ three-valued
\begin{equation*}
    \rho  (\mathbf{r },t)= \Psi ^\dagger(\mathbf{r },t)\Psi(\mathbf{r },t) \qquad
    \mathbf{j}(\mathbf{r },t)=\Psi ^\dagger (\mathbf{r },t)\hat{\boldsymbol{\alpha } }\Psi(\mathbf{r },t)
\end{equation*}
which does not suffer from non-conserved probability density
\begin{equation*}
    \partial_t \int d^3r\;\rho =0
\end{equation*}

\subsubsection{Derivation: Dirac approach.}
He considered a wave equation with only first derivatives with respect to time and space and introduced the classical energy
\begin{equation*}
    E = c\boldsymbol{\hat{\alpha}}\cdot\mathbf{p} + \hat{\beta}m c^2
\end{equation*}
such that squaring it one recovers the relativistic energy
\begin{equation*}
    E ^{2 } =  \mathbf{p }^2c^2 +m^2c^4
\end{equation*}
Since spatial translations are generated by momentum and time translations are generated by the Hamiltonian, we introduce quantization prescription
\begin{equation*}
    \mathbf{p} \to -i\hbar\nabla \qquad
    H \to i\hbar \;\partial_t
\end{equation*}
The classical energy now reads
\begin{equation*}
    i\hbar \;\partial_t \Psi(\mathbf{r},t)=\left(-i\hbar c\hat{\boldsymbol{\alpha}}\cdot\nabla + \hat{\beta}m c^{2}\right)\Psi(\mathbf{r},t)
\end{equation*}
Notice that the wavefunction $\Psi(\mathbf{r },t)$, has four components in the abstract space of Dirac matrices, as such
\begin{equation*}
    i\hbar\partial_t
    \begin{bmatrix}
        \psi_1(\mathbf{r},t) \\
        \psi_2(\mathbf{r},t) \\
        \psi_3(\mathbf{r},t) \\
        \psi_4(\mathbf{r},t)
    \end{bmatrix}
    =
    \hat{H}
    \begin{bmatrix}
        \psi_1(\mathbf{r},t) \\
        \psi_2(\mathbf{r},t) \\
        \psi_3(\mathbf{r},t) \\
        \psi_4(\mathbf{r},t)
    \end{bmatrix}
\end{equation*}

The first term of the Hamiltonian has the elements of
\begin{align*}
    \boldsymbol{\sigma}\cdot \hat{\mathbf p}
     & = -i\hbar\left( \sigma_x \partial_x + \sigma_y \partial_y + \sigma_z \partial_z \right) \\
     & = -i\hbar \left(
    \begin{bmatrix}
            0          & \partial_x \\
            \partial_x & 0
        \end{bmatrix}
    +
    \begin{bmatrix}
            0           & -i\partial_y \\
            i\partial_y & 0
        \end{bmatrix}
    +
    \begin{bmatrix}
            \partial_z & 0           \\
            0          & -\partial_z
        \end{bmatrix}
    \right)                                                                                    \\
     & = -i\hbar
    \begin{bmatrix}
        \partial_z               & \partial_x - i\partial_y \\
        \partial_x + i\partial_y & -\partial_z
    \end{bmatrix}
\end{align*}
Along the second term being trivial, the Hamiltonian explicitly written as
\begin{equation*}
    \begin{bmatrix}
        mc^2                                             & 0                                                & -i\hbar c\;\partial_z                            & -i \hbar c \left( \partial_x-i \;\partial_y \right) \\
        0                                                & mc^2                                             & -i\hbar c\left(\partial_x + i\;\partial_y\right) & i\hbar c \partial_z                                 \\
        -i\hbar c\;\partial_z                            & -i\hbar c\left(\partial_x - i\;\partial_y\right) & -mc^2                                            & 0                                                   \\
        -i\hbar c\left(\partial_x + i\;\partial_y\right) & i\hbar c\;\partial_z                             & 0                                                & -mc^2
    \end{bmatrix}
\end{equation*}

\subsubsection{Derivation: continuity equation.}
Start from the Dirac equation and its conjugate
\begin{gather*}
    \partial_t \Psi(\mathbf{r},t)=-c\hat{\boldsymbol{\alpha}}\cdot\nabla \Psi(\mathbf{r},t)
    + (i m c^{2}/\hbar )\hat{\beta} \Psi(\mathbf{r},t) \\
    \partial_t \Psi ^\dagger(\mathbf{r},t)=-c [\nabla \Psi ^\dagger(\mathbf{r},t)] \cdot \hat{\boldsymbol{\alpha}}
    -(i m c^{2}/\hbar )\Psi ^\dagger(\mathbf{r},t) \hat{\beta}
\end{gather*}
Left multiply the Dirac equation with the Hermitian conjugate of the wavefunction
\begin{equation*}
    \Psi ^\dagger \partial_t \Psi=c\Psi ^\dagger\hat{\boldsymbol{\alpha}}\cdot\nabla \Psi
    + (i m c^{2}/\hbar ) \Psi ^\dagger \hat{\beta} \Psi
\end{equation*}
the right multiply the Hermitian conjugate of the Dirac equation with the wavefunction
\begin{equation*}
    (\partial_t \Psi ^\dagger) \Psi=c [\nabla \Psi ^\dagger] \cdot \hat{\boldsymbol{\alpha}}\Psi
    -(i m c^{2}/\hbar )\Psi ^\dagger \hat{\beta} \Psi
\end{equation*}
and add both
\begin{equation*}
    \partial_t \Psi ^\dagger\Psi
    = -c \left[     \Psi ^\dagger\hat{\boldsymbol{\alpha}}\cdot\nabla \Psi + (\nabla \Psi ^\dagger) \cdot \hat{\boldsymbol{\alpha}}\Psi \right]
\end{equation*}

Noting that $\hat{\boldsymbol{\alpha}}$ is constant in space such that $\nabla \hat{\boldsymbol{\alpha}} = 0$, can write
\begin{align*}
    \nabla \cdot (\Psi ^\dagger \hat{\boldsymbol{\alpha } }\Psi)
     & = \sum_{i } \partial_i\big(\Psi^\dagger\hat{\alpha}_i\Psi\big)                                                                                      \\
     & =  \sum_{i } (\partial_i\Psi^\dagger) \hat{\alpha}_i \Psi + \Psi^\dagger(\partial_i\hat{\alpha}_i)\Psi + \Psi^\dagger\hat{\alpha}_i(\partial_i\Psi) \\
     & = \sum_i\big[(\partial_i\Psi^\dagger)\hat{\alpha}_i \Psi + \Psi^\dagger\hat{\alpha}_i(\partial_i\Psi)\big]                                          \\
    \nabla \cdot (\Psi ^\dagger \hat{\boldsymbol{\alpha } }\Psi)
     & = (\nabla\Psi^\dagger)\cdot\boldsymbol{\hat{\alpha}}\Psi + \Psi^\dagger\big(\boldsymbol{\hat{\alpha}}\cdot\nabla\Psi\big)
\end{align*}
So we obtain the continuity equation
\begin{equation*}
    \partial_t \rho + c \nabla \cdot \mathbf{j}=0
\end{equation*}
where
\begin{equation*}
    \rho  = \Psi ^\dagger\Psi \qquad
    \mathbf{j}=\Psi ^\dagger \hat{\boldsymbol{\alpha } }\Psi
\end{equation*}

Now integrate over all space
\begin{equation*}
    \int d^3r\;\partial_t \rho + \int d^3r\;\nabla \cdot  \mathbf{j} = 0
\end{equation*}
Using divergence theorem
\begin{equation*}
    \partial_t \int d^3r\;\rho + \oint d \mathbf{a}\cdot  \mathbf{j} = 0
\end{equation*}
Applying vanishing current at infinity, we obtain the conserved probability density
\begin{equation*}
    \partial_t \int d^3r\;\rho =0
\end{equation*}

\subsection{Pauli Equation}
The origin of spin can be traced back in the Pauli equation.
The standard form is
\begin{equation*}
    i\hbar \; \partial_t\psi
    =  \left[\frac{1}{2m}(\mathbf{{\hat{p}}}-q\mathbf{A}) ^{2}-\frac{q \hbar}{2m }\mathbf{B }\cdot  \boldsymbol{\sigma} +q\phi\right]\psi
\end{equation*}
Explicitly, the spin appear as
\begin{equation*}
    i\hbar \; \partial_t\psi
    =  \left[\frac{1}{2m}(\mathbf{{\hat{p}}}-q\mathbf{A}) ^{2}-\frac{q }{m }\mathbf{B }\cdot \mathbf{S} +q\phi\right]\psi
\end{equation*}

\subsubsection{Derivation.}
Suppose the relativistic particle has the electric charge $q$, by using the Gauge-invariant substitution
\begin{equation*}
    E \to i\hbar\;\partial_t-q\phi(\mathbf r,t)  \qquad \mathbf{p} \to -i\hbar\nabla - q\mathbf{A}(\mathbf r,t)
\end{equation*}
We obtain the Dirac equation in the form
\begin{equation*}
    i\hbar\;\partial_t\Psi(\mathbf r,t) =
    \Big(c\hat{\boldsymbol\alpha}\cdot\big(\hat{\mathbf p}-q\mathbf{A}(\mathbf r,t)\big)+\beta mc^2+q\phi(\mathbf r,t)\Big)\Psi(\mathbf{ r},t)
\end{equation*}
The spinor then is set as
\begin{equation*}
    \Psi(\mathbf r,t)
    \to e^{-imc^2t/\hbar}
    \begin{bmatrix}
        \psi(\mathbf r,t) \\
        \chi(\mathbf r,t)
    \end{bmatrix}
\end{equation*}
so that we obtain
\begin{align*}
    (    i\hbar \; \partial_t- mc^2)
    \begin{bmatrix}
        \psi \\
        \chi
    \end{bmatrix}
     & =
    \begin{bmatrix}
        q\phi+mc^2                                             & c\boldsymbol{\sigma}\cdot(\hat{\mathbf p}-q\mathbf{A}) \\
        c\boldsymbol{\sigma}\cdot(\hat{\mathbf p}-q\mathbf{A}) & q\phi-mc^2
    \end{bmatrix}
    \begin{bmatrix}
        \psi \\
        \chi
    \end{bmatrix}
    \\
    i\hbar \; \partial_t
    \begin{bmatrix}
        \psi \\
        \chi
    \end{bmatrix}
     & =
    \begin{bmatrix}
        q\phi                                                  & c\boldsymbol{\sigma}\cdot(\hat{\mathbf p}-q\mathbf{A}) \\
        c\boldsymbol{\sigma}\cdot(\hat{\mathbf p}-q\mathbf{A}) & q\phi-2mc^2
    \end{bmatrix}
    \begin{bmatrix}
        \psi \\
        \chi
    \end{bmatrix}
\end{align*}

Consider non-relativistic approximation
\begin{equation*}
    \bigl(i\hbar \; \partial_t - q\phi + 2mc^{2}\bigr)\chi \approx 2mc^{2}\chi
\end{equation*}
The term $i\hbar \; \partial_t - q\phi$ represent kinetic operator in this case, since it is equal to total energy $i\hbar \; \partial_t$ minus the potential $q\phi$ operator.
The non-relativistic condition reads $v \ll c$, which this approximation satisfies by $T \sim mv^2/2$.
Under this approximation, the equation reads
\begin{equation*}
    i\hbar \; \partial_t
    \begin{bmatrix}
        \psi \\
        0
    \end{bmatrix}
    =
    \begin{bmatrix}
        q\phi                                                  & c\boldsymbol{\sigma}\cdot(\hat{\mathbf p}-q\mathbf{A}) \\
        c\boldsymbol{\sigma}\cdot(\hat{\mathbf p}-q\mathbf{A}) & -2mc^2
    \end{bmatrix}
    \begin{bmatrix}
        \psi \\
        \chi
    \end{bmatrix}
\end{equation*}

We have then the lower equation
\begin{equation*}
    \chi = \frac{\hat{\boldsymbol{\sigma}}\cdot(\hat{\mathbf{p}}-qA)}{2mc}\psi
\end{equation*}
Inserting this expression in the upper equation
\begin{equation*}
    i\hbar \; \partial_t\psi
    = \left(\frac{\big[\hat{\sigma}\cdot(\hat{\mathbf{p}}-qA)\big]^2}{2m}+q\phi\right)\psi
\end{equation*}
Using
\begin{equation*}
    (\mathbf{a} \cdot \boldsymbol{\sigma})(\mathbf{b} \cdot \boldsymbol{\sigma}) = \mathbf{a} \cdot \mathbf{b} + i(\mathbf{a} \times \mathbf{b}) \cdot \boldsymbol{\sigma}
\end{equation*}
we can evaluate
\begin{align*}
    (\boldsymbol{\sigma}\cdot(\mathbf{{\hat{p}}}-q\mathbf{A}))^2
     & = (\mathbf{{\hat{p}}}-q\mathbf{A}) ^{2}+ i\big((\mathbf{{\hat{p}}}-q\mathbf{A})\times(\mathbf{{\hat{p}}}-q\mathbf{A})\big)\cdot\boldsymbol{\sigma}    \\
     & = (\mathbf{{\hat{p}}}-q\mathbf{A}) ^{2}+ i\big(    \mathbf{{\hat{p}}}\times\mathbf{{\hat{p}}} - q\mathbf{{\hat{p}}}\times\mathbf{A}                   \\
     & \qquad- q\mathbf{A}\times\mathbf{{\hat{p}}} + q^{2}\mathbf{A}\times\mathbf{A}\big)\cdot\boldsymbol{\sigma}                                            \\
     & = (\mathbf{{\hat{p}}}-q\mathbf{A}) ^{2}+ iq\big(\mathbf{{\hat{p}}}\times\mathbf{A} + \mathbf{A}\times\mathbf{{\hat{p}}} \big)\cdot\boldsymbol{\sigma} \\
    (\boldsymbol{\sigma}\cdot(\mathbf{{\hat{p}}}-q\mathbf{A}))^2
     & = (\mathbf{{\hat{p}}}-q\mathbf{A}) ^{2}+ q \hbar \mathbf{B}\cdot\boldsymbol{\sigma}
\end{align*}
The anticommutator is evaluated on test function as
\begin{align*}
    \left\{ \mathbf{\hat{p },\mathbf{A }} \right\} \psi
     & =-i\hbar\nabla\times\mathbf{A}\psi - i \hbar\mathbf{A}\times\nabla\psi                                                  \\
     & = -i\hbar\big[(\nabla\times\mathbf{A}) \psi + (\nabla \psi) \times  \mathbf{A}+\mathbf{A}\times(\nabla\psi)\big]        \\
     & = -i\hbar\big[(\nabla\times\mathbf{A}) \psi + (\nabla \psi) \times  \mathbf{A} - (\nabla \psi) \times  \mathbf{A} \big] \\
    \left\{ \mathbf{\hat{p },\mathbf{A }} \right\} \psi
     & = -i\hbar(\nabla\times\mathbf{A}) \psi = -i\hbar \mathbf{B} \psi
\end{align*}
The spin naturally occur by design
\begin{align*}
    i\hbar \; \partial_t\psi
    =  \left[\frac{1}{2m}(\mathbf{{\hat{p}}}-q\mathbf{A}) ^{2}-\frac{q \hbar}{2m }\mathbf{B }\cdot  \boldsymbol{\sigma} +q\phi\right]\psi \\
    i\hbar \; \partial_t\psi
    =  \left[\frac{1}{2m}(\mathbf{{\hat{p}}}-q\mathbf{A}) ^{2}-\frac{q }{m }\mathbf{B }\cdot \mathbf{S} +q\phi\right]\psi
\end{align*}
It is an unavoidable consequence of the Dirac algebra structure once the non-relativistic limit is taken.

\end{document}